\section{\texorpdfstring{Proofs for Sections~\ref{sec:multistate-model}
and~\ref{sec:simplex-thresholds}}{Proofs for Sections 2 and 3}}\label{app:thresholds}

\subsection{Facts about Bessel series}

\begin{lemma}[the Bessel law]\label{lem:bessel-law}
For $x\ge0$ let $J$ be a random variable with
$\Prob(J=m)=x^m/[m!(m+1)!\,\Phi(x)]$ for $m\ge0$ (so $J=0$ when $x=0$), and put
$\psi(x)=\E J$. Then
\begin{enumerate}[(i)]
\item $\E[J(J+1)]=x$,
\item $\E[J(J-1)J(J+1)]=x^2$,
\item $\psi(x)=x\Phi'(x)/\Phi(x)$, and $\psi(x)=\sqrt x\,I_2(2\sqrt x)/I_1(2\sqrt x)$
for $x>0$,
\item $x\psi'(x)=\Var J$, so $0\le\psi'(x)\le1$,
\item $0\le\psi(x)\le\min(x/2,\sqrt x)$,
\item $0\le(\log\Phi)'(x)=\psi(x)/x\le\frac12$ for $x>0$.
\end{enumerate}
\end{lemma}

\begin{proof}
(i) We have $\sum_{m\ge1}m(m+1)x^m/[m!(m+1)!]=\sum_{m\ge1}x^m/[(m-1)!m!]=x\Phi(x)$.
(ii) For $m\ge2$, $m(m-1)/m!=1/(m-2)!$ and $m(m+1)/(m+1)!=1/(m-1)!$, so the sum is
$\sum_{m\ge2}x^m/[(m-2)!(m-1)!]=x^2\Phi(x)$.
(iii) The first form is the definition of the mean. Termwise,
$\Phi'(x)=\sum_{i\ge0}x^i/[i!(i+2)!]=I_2(2\sqrt x)/x$ by the series of
$I_2$~\cite[Eq.~10.25.2]{dlmf}, which gives the second.
(iv) Since $\log\Prob(J=m)=m\log x-\log\Phi(x)$ plus a term free of $x$, we
have $x\frac d{dx}\log\Prob(J=m)=m-\psi(x)$. So $x\psi'(x)=\E[J(J-\psi)]=\Var J$,
and $\Var J\le\E J^2\le\E[J(J+1)]=x$.
(v) For integers $J\ge0$ we have $J\le J(J+1)/2$, so $\psi\le x/2$ by (i). Also
$\psi^2\le\E J^2\le x$.
(vi) follows from (iii) and (v).
\end{proof}

\begin{lemma}[clipped product deficit]\label{lem:clipped-product}
For any finite list of nonnegative real numbers $t_1,\ldots,t_m$, with
$(x)_+=\max(x,0)$,
\[
0\le1-\prod_{j=1}^m(1-t_j)_+\le\sum_{j=1}^mt_j.
\]
\end{lemma}

\begin{proof}
If some $t_j\ge1$, the product is zero and the sum is at least one. Otherwise all
factors lie in $[0,1]$, and induction follows from
$1-(1-a)(1-b)=a+b-ab\le a+b$. The empty product gives equality.
\end{proof}

\begin{lemma}[Gaussian averages of the factors $\Phi$]\label{lem:gaussian-phi-average}
Fix $q\ge0$, an integer $\ell\ge0$ and $X\ge0$. Let $w_0\ge1$ and
$c_1,\ldots,c_\ell\ge0$ with $\xi_j=c_jw_0^2\le X$, and put
$F(w)=w^q\prod_{j=1}^\ell\Phi(c_jw^2)$. Let $J_1,\ldots,J_\ell$ be independent,
with $J_j$ distributed as $J$ in Lemma~\ref{lem:bessel-law} with $x=\xi_j$. Put
$|J|=J_1+\cdots+J_\ell$ and $\mu_q=\E[(q+2|J|)(q+2|J|-1)]$. Then
\[
\frac1{\sqrt\pi}\int_{-w_0/2}^{w_0/2}e^{-y^2}F(w_0+y)\,dy
=F(w_0)\Bigl(1+\frac{\mu_q}{4w_0^2}+O(w_0^{-4})\Bigr),
\]
and $F(w_0+y)=O(F(w_0))$ for $|y|\le w_0/2$. The constants depend only on $q$,
$\ell$ and $X$. Moreover
\[
\mu_q=q^2-q+4\sum_j\xi_j+(4q-6)\sum_j\psi(\xi_j)
+8\sum_{i<j}\psi(\xi_i)\psi(\xi_j).
\]
\end{lemma}

\begin{proof}
Expanding each $\Phi$ gives $F(w)=\sum_m\kappa_mw^{q+2|m|}$, where
$m=(m_1,\ldots,m_\ell)$ runs over vectors of nonnegative integers, $|m|=\sum_jm_j$
and $\kappa_m=\prod_jc_j^{m_j}/[m_j!(m_j+1)!]\ge0$. The weights
$\kappa_mw_0^{q+2|m|}/F(w_0)$ are the joint law of $(J_1,\ldots,J_\ell)$. Hence
$F^{(i)}(w_0)=F(w_0)\,\E[(q+2|J|)_i]\,w_0^{-i}$, where
$(y)_i=y(y-1)\cdots(y-i+1)$. In particular $F''(w_0)=\mu_qw_0^{-2}F(w_0)$.

Let $w=w_0+y$ with $|y|\le w_0/2$. For every $\eta\ge0$ we have
$w^{\eta-4}\le16(3/2)^\eta w_0^{\eta-4}$ and $|(\eta)_4|\le\eta^4+1$. Applying this
with $\eta=q+2|m|$ termwise gives
\[
|F''''(w)|\le16\Bigl(\frac32\Bigr)^qw_0^{q-4}
\prod_j\Phi\Bigl(\frac{9\xi_j}4\Bigr)\,
\tilde\E\bigl[(q+2|\tilde J|)^4+1\bigr],
\]
where the $\tilde J_j$ are independent with the parameters $9\xi_j/4$ in place of
$\xi_j$. By Lemma~\ref{lem:bessel-law}(i) and (ii),
$\E J^4=x^2+\E J^2\le x^2+x$. Hence $\tilde\E[(q+2|\tilde J|)^4]$ is bounded in
terms of $q$, $\ell$ and $X$. Since $\Phi\ge1$ we have $F(w_0)\ge w_0^q$, so
$F''''(w)=O(w_0^{-4}F(w_0))$. The same argument without derivatives bounds
$F(w_0+y)$.

Now expand $F(w_0+y)$ to third order in $y$ with the Lagrange remainder. The odd
terms integrate to zero over the symmetric interval. The integrals of $e^{-y^2}$
and $y^2e^{-y^2}$ over $|y|\le w_0/2$ differ from $\sqrt\pi$ and $\sqrt\pi/2$ by
$O(w_0e^{-w_0^2/4})$, the number $\mu_q$ is bounded, and the remainder term is
$O(w_0^{-4}F(w_0))$. This proves the expansion.

For the last formula put $\psi_j=\psi(\xi_j)=\E J_j$. Then
$\E J_j^2=\xi_j-\psi_j$ by Lemma~\ref{lem:bessel-law}(i), and independence gives
$\E[(q+2|J|)^2]=q^2+4q\sum_j\psi_j+4\sum_j(\xi_j-\psi_j)+8\sum_{i<j}\psi_i\psi_j$.
Subtracting $\E[q+2|J|]=q+2\sum_j\psi_j$ proves it.
\end{proof}

\subsection{Proofs for Section~\ref{sec:simplex-thresholds}}

\begin{proof}[Proof of the exact assertions in Theorem~\ref{thm:simplex-crossing}]
Let $n$ be a balanced bin with $k$ positive coordinates, so that
$C_{N,k}/V_N=S_n(u)$ with $u=(1-p)/((d-1)p)$. To count $W_n(j)$, choose a
sequence of $j+1$ run colors in which successive colors differ. If color $i$
appears $m_i$ times, its run lengths can be chosen in $\binom{n_i-1}{m_i-1}$
ways. So every coefficient of $S_n$ is nonnegative, and the first nonzero one is
$W_n(k-1)=k!$. Hence $S_n$ increases strictly from $0$ to $\infty$ on $u>0$ and
equals $1$ exactly once. As $p$ increases over $(0,1)$, $u$ decreases over
$(0,\infty)$, which gives the unique $p_{N,k}$ and the 2 inequalities.

Choose the balanced bins for increasing $N$ by adding $1$ to one of the smallest
coordinates. Each coefficient $W_n(j)$ is then nondecreasing. To see this, put an
extra copy of the incremented color just before its first occurrence. This keeps
the number of changes, and different words give different words. The coefficient
of $u^k$ is
\[
W_n(k)=\frac{k!(k-1)}2\sum_i(n_i-1)=\frac{k!(k-1)}2(N-k).
\]
Indeed, exactly one run color appears twice. For each choice of it there are
$(k+1)!/2-k!=k!(k-1)/2$ sequences of run colors in which successive colors
differ, and $n_i-1$ ways to split its $n_i$ visits into 2 runs. This coefficient
increases strictly with $N$, so the root of $S_n(u)=1$ decreases strictly with
$N$, and $p=1/(1+(d-1)u)$ increases. If $N=k$, every coordinate is $1$ and
$S_n(u)=k!u^{k-1}$, which gives $p_{k,k}$. Since the root decreases with $N$,
every root has $u\le(k!)^{-1/(k-1)}<1$.
\end{proof}

The asymptotic assertions rest on a Laplace expansion of the positive
sum~\eqref{eq:simplex-positive-sum}. It is the case with no rare counts of the
crossover theorem in Appendix~\ref{app:boundary}.

\begin{lemma}[no rare counts]\label{lem:no-rare}
Fix $k\ge2$, $\varepsilon>0$ and $0<c<C$. Let $n$ be a positive count vector
with $k$ coordinates and total $N$, put $\beta=n/N$, and assume that every
$\beta_i\ge\varepsilon$. Put $A=\sum_i\sqrt{\beta_i}$, $g=A^2-1$ and $s=k-1$,
let $D(\beta)$ be as in Theorem~\ref{thm:simplex-interior}, and put
\[
\gamma_0(\beta)=\frac{k(k+2)}{16A^2}-\frac3{16A}\sum_{i=1}^k\beta_i^{-1/2}.
\]
Then, uniformly for these $n$ and $cL\le z\le CL$,
\begin{equation}\label{eq:no-rare}
S_n(z/N)=D(\beta)\frac{z^{s/2}e^{gz}}{N^s}\,e^{\gamma_0(\beta)/z}
\Bigl(1+O\Bigl(\frac1{L^2}+\frac{L^2}N\Bigr)\Bigr).
\end{equation}
\end{lemma}

\begin{proof}
This is Theorem~\ref{thm:simplex-boundary-crossover} with $\ell=0$ and $k_1=k$.
Then $M=N$, $\eta=s/2$, the constant $D_{k_1}(\beta)$ is $D(\beta)$, there are
no factors $\Phi$, and $\mathcal E(z)=\gamma_0/z$ with $\gamma_0$ as above,
since $(2\eta+1)(2\eta+3)=k(k+2)$. The proof of that theorem uses only
Theorem~\ref{thm:uniform-simplex-majorant}, the expansion of $I_1$ and
Lemma~\ref{lem:gaussian-phi-average}, and nothing from
Section~\ref{sec:simplex-thresholds}.
\end{proof}

\begin{proof}[Proof of the asymptotic assertions in
Theorems~\ref{thm:simplex-window} and~\ref{thm:simplex-crossing}]
Put $s=k-1$ and let $n$ be balanced, so that $\beta_i=n_i/N=1/k+O(1/N)$. At the
center $A^2-1=s$, $D(\beta)=D_k$ and $\gamma_0=(k+2)/16-3k/16=-s/8$. On the
plane $\sum_i\beta_i=1$ the function $A^2$ has a critical point at the center,
so $A^2-1=s+O(N^{-2})$, while $D(\beta)$ and $\gamma_0(\beta)$ change by
$O(1/N)$. Since $z=O(L)$, Lemma~\ref{lem:no-rare} gives
\eqref{eq:simplex-uniform-ratio}.

Now let $cL\le T\le CL$ and $u=r/p$, so that $z=T/p$ lies in $[cL,2CL]$ for large
$N$ and $z-T=(d-1)Tz/N=O(L^2/N)$. Take logarithms
in~\eqref{eq:simplex-uniform-ratio} and replace $z$ by $T$. This changes each term
by $O(T^2/N)$, and $\log D_k=-sa_k$, so we get (i). For (ii),
$\log V_N=(N-1)\log(1-(d-1)T/N)-\log d=-(d-1)T-\log d+O(T^2/N)$. With $T=\tau L$
and (i) we get $\log C_{N,k}=-(d-1)\tau L+(k-1)(\tau-1)L+O(\log L)$, which is (ii),
and the case $k=1$ is the formula for $\log V_N$. For (iii),
$T=L-\frac12\log L+t$ gives $\Psi_N(T)=t+O(\log L/L)$ uniformly on compact $t$
sets.

For the crossing, (i) at $T=(1-\varepsilon)L$ and at $T=(1+\varepsilon)L$ shows
that the ratio tends to $0$ and to $\infty$. Since $S_n(u)$ increases with $u$
and $u$ increases with $T$, the root has $z_{N,k}/L\to1$. At the root the left
side of (i) is $0$, and dividing by $s$ gives the implicit equation, since
$T^2/N=o(L^{-2})$. For the explicit form put
$z_*=L-\frac12\log L+a_k+e_L/L$ with $e_L=\frac14\log L-\frac12a_k+\frac18$.
Taylor expansion gives $z_*+\frac12\log z_*-1/(8z_*)=L+a_k+O((\log L)^2/L^2)$.
The derivative of $z+\frac12\log z-1/(8z)$ is at least $1$ for $z>0$, so
$z_{N,k}-z_*=O((\log L)^2/L^2)$, which is~\eqref{eq:simplex-root-expansion}.

For (iv), $a_k=\frac12\log(4\pi)-f(k)$ with $f(x)=(x+\frac12)\log x/(x-1)$. For
$x\ge2$ the sign of $f'(x)$ is the sign of $\Delta(x)=2x^2-x-1-3x\log x$. Here
$\Delta(2)=5-6\log2>0$, and $\Delta'(x)=4x-4-3\log x>0$ for $x\ge2$, since
$\Delta'(2)>0$ and $\Delta''(x)=4-3/x>0$. So $f$ increases strictly and $a_k$
decreases strictly.
\end{proof}

\begin{proof}[Proof of Theorem~\ref{thm:simplex-interior}]
The polynomial $S_n$ has nonnegative coefficients and lowest term $k!u^{k-1}$,
so $S_n(u)=1$ has a unique positive root. Lemma~\ref{lem:no-rare} with
$\beta=\alpha$ gives~\eqref{eq:simplex-interior-ratio}, since $\gamma_0(\alpha)$
is bounded when every $\alpha_i\ge\varepsilon$.

On the compact set of compositions, $g\ge2\varepsilon$ and $D(\alpha)$ has
positive lower and finite upper bounds. Testing the ratio at suitable constant
multiples of $\log N$ brackets the root uniformly. Taking logarithms there gives
\[
gz+\frac s2\log z=sL-\log D(\alpha)+O_{k,\varepsilon}(z^{-1}+z^2/N).
\]
First $z=(s/g)L+O_{k,\varepsilon}(\log L)$, and one substitution then
gives~\eqref{eq:simplex-interior-root}. The derivative of the left side is at
least $g$, so the error is preserved on inversion. Finally, the functions of
$\alpha$ in the formula have bounded derivatives on a slightly smaller compact
set. So moving $\alpha$ by $O(1/N)$ changes the expression by $O(L/N)$, which the
error term absorbs.
\end{proof}

\begin{proof}[Proof of Lemma~\ref{lem:occupation-covariance}]
Note that $(\mathcal Lx)_i=\sum_jc_{ij}(x_i-x_j)$ and $\mathcal L\mathbf1=0$. Take
$\theta$ with $\theta_1=0$ and put $\zeta=\operatorname{diag}(\pi)f$. Then
$\sum_i\zeta_i=0$, $\mathcal L\chi=\zeta$ and the form is $2\zeta^\top\chi$. The
vector $\tilde\chi=\chi-\chi_1\mathbf1$ also solves $\mathcal L\tilde\chi=\zeta$,
with $\zeta^\top\tilde\chi=\zeta^\top\chi$ and $\tilde\chi_1=0$. So rows $2$ to
$k$ give $\mathcal L_{(1)}\tilde\chi'=\zeta'$, where primes delete the first
coordinate, and
the form is $2\zeta'^\top\mathcal L_{(1)}^{-1}\zeta'$. As $\theta_1=0$, we have
$\zeta'=(D'-\pi'\pi'^\top)\theta'$ with $D'=\operatorname{diag}(\pi_2,\ldots,\pi_k)$.
Hence $\Sigma=2(D'-\pi'\pi'^\top)\mathcal L_{(1)}^{-1}(D'-\pi'\pi'^\top)$. The
matrix determinant lemma gives
$\det(D'-\pi'\pi'^\top)=\prod_{i\ge2}\pi_i\,(1-\sum_{i\ge2}\pi_i)=\prod_i\pi_i$,
and the matrix-tree theorem~\cite{bollobas1998} gives
$\det\mathcal L_{(1)}=\tree(c)>0$.
\end{proof}

\begin{proof}[Proof of Proposition~\ref{prop:kirchhoff}]
(1) For $i\ne j$ we have $\alpha_iG^\alpha_{ij}=h_ih_j$, which is symmetric, so
$\alpha$ is reversible and hence stationary. Also $(Q_kh)_i=A-kh_i$, so
$\theta_{\alpha,i}=k-A/h_i$, and $Q_k+\operatorname{diag}(\theta_\alpha)$ has
off-diagonal entries $1$ and diagonal entries $1-A/h_i$. It is symmetric and
kills $h>0$, so $0$ is its Perron root. Conjugating by $\operatorname{diag}(h)$
gives the rates $h_j/h_i$. If $G^\alpha$ were a Doob transform of $Q_k$, then $h$
would be an eigenvector of $Q_k$, so $A-kh_i=\lambda h_i$ for all $i$ and $\alpha$
would be uniform.
(2) is Lemma~\ref{lem:weighted-cayley} with $x_i=h_i$.
(3) By (2), $\sqrt{\tree(\alpha)}/\prod_i\alpha_i=(\prod_i\alpha_i)^{-3/4}A^{k/2-1}$,
which gives the first form. Lemma~\ref{lem:occupation-covariance} with $\pi=\alpha$
gives $\det(2\pi\Sigma_\alpha)=(4\pi)^{k-1}(\prod_i\alpha_i)^2/\tree(\alpha)$, which
gives the second.
(4) At the center $\mathcal L_{(1)}=I-J/k$ and $D'-\alpha'\alpha'^\top=(I-J/k)/k$
in the formula $\Sigma=2(D'-\alpha'\alpha'^\top)\mathcal L_{(1)}^{-1}
(D'-\alpha'\alpha'^\top)$ from the proof of
Lemma~\ref{lem:occupation-covariance}. This gives $\Sigma_k$, and the rest is
substitution.
(5) is the same computation with $k_1$ and $\beta$.
\end{proof}

\begin{proof}[Proof of Proposition~\ref{prop:mass-balance}(2) and (3)]
Let $u=z/N$ with $cL\le z\le CL$, and put $s=k-1$. In (1) the term $j=k$ is
$k(1+(k-1)z/N)^{N-1}=k\,e^{(k-1)z+O(z^2/N)}$, and each term with $j<k$ is smaller
by a factor $O(e^{-z})$. This gives the first expression.

For the second, write $\alpha'=(\alpha_2,\ldots,\alpha_k)$ and $g(\alpha)=A^2-1$ as
in Theorem~\ref{thm:simplex-interior}. In full coordinates the Hessian of $A^2$ at
the center is $(k/2)J-(k^2/2)I$. Pulling it back along $w=(-\sum_iv_i,v)$ gives
$-\nabla^2g=(k^2/2)(I+J)$ in the coordinates $\alpha'$, and its inverse is
$(2/k^2)(I-J/k)=\Sigma_k$. Fix a small $\eta>0$ and let $K_\eta$ be the set of
$\alpha$ with $|\alpha_i-1/k|\le\eta$ for all $i$. On $K_\eta$,
Theorem~\ref{thm:simplex-interior} gives
$S_n=D(\alpha)z^{s/2}e^{g(\alpha)z}N^{-s}(1+O(z^{-1}+z^2/N))$ uniformly, with
$\alpha=n/N$. The points $\alpha'$ form a lattice with cells of volume $N^{-s}$,
and on each cell $g$ changes by $O(1/N)$ and $D$ by a factor $1+O(1/N)$. So the
sum over $n/N\in K_\eta$ is $N^s$ times the integral of the same expression over
$K_\eta$, times $1+O(z/N)$, up to the cells near the boundary of $K_\eta$. There
$g\le k-1-c_\eta$ for some $c_\eta>0$, so these cells are negligible. The function
$g$ has its strict maximum $k-1$ at the center, and $D$ is continuous there with
$D=D_k$. So Laplace's method gives
\[
\int_{K_\eta}D(\alpha)e^{g(\alpha)z}\,d\alpha'
=D_ke^{(k-1)z}\Bigl(\frac{2\pi}z\Bigr)^{s/2}(\det\Sigma_k)^{1/2}(1+o(1)).
\]
Multiplied by $z^{s/2}$ this is the second expression, provided the bins off
$K_\eta$ are negligible. For them, Theorem~\ref{thm:uniform-simplex-majorant} and
$\Phi(x)\le e^{2\sqrt x}$ give, with $Y=Nv$ and $A_n=\sum_i\sqrt{n_i/N}$,
\[
S_n(u)\le(1-u)^{N-1}v^s\int_0^\infty e^{-t+2A_n\sqrt{Yt}}\,t^k\,dt
\le C\,\frac{(1+z)^{2k}}{N^s}\,e^{(A_n^2-1)z+O(z^2/N)} .
\]
Here $A_n^2\le k-c'_\eta$ for some $c'_\eta>0$, since $A^2$ is continuous on the
closed simplex and equals $k$ only at the center. There are at most $N^s$ bins,
so they contribute $O(z^{2k}e^{(k-1-c'_\eta)z})=o(e^{(k-1)z})$. This proves (2).

For (3), equate the 2 expressions in (2) and let $N\to\infty$. This gives
$k=D_k(2\pi)^{s/2}(\det\Sigma_k)^{1/2}$, that is, $D_k=k\det(2\pi\Sigma_k)^{-1/2}$.
As a check, $\det\bigl((k^2/2)(I+J)\bigr)=(k^2/2)^sk$, which gives
$D_k=k^{k+1/2}(4\pi)^{-s/2}$ again.
\end{proof}

\begin{proof}[Proof of Theorem~\ref{thm:simplex-balance}]
At $p=1/d$ this follows from the multinomial formula. For $1/d<p<1$ we have
$\sigma>0$ and $v=r/\sigma>0$. Writing $M!=\int_0^\infty e^{-t}t^M\,dt$ in the
uniform case of Proposition~\ref{prop:dpmf}(a) gives
\begin{equation}\label{eq:simplex-balance-integral}
P_N(n)=\frac{\sigma^{N-1}}{dv}\int_0^\infty e^{-t}
\prod_iB_{n_i}(vt)\,dt,\qquad
B_n(y)=\sum_{m=1}^n\binom{n-1}{m-1}\frac{y^m}{m!},
\end{equation}
where the product runs over the positive coordinates. For each $y>0$ the sequence
$B_n(y)$, $n\ge1$, is strictly log-concave, and we give a direct proof. Write
\[
\frac{B_n(y)}y=b_{n-1},\qquad
b_j=\sum_{i=0}^j\binom ji a_i,\quad
a_i=\frac{y^i}{(i+1)!},\qquad
c_j=\sum_{i=0}^j\binom ji a_{i+1}.
\]
Pascal's identity gives $b_{j+1}=b_j+c_j$. The ratio $c_j/b_j$ is the mean of the
strictly decreasing sequence $a_{i+1}/a_i=y/(i+2)$ under weights proportional to
$\binom jia_i$. When $j$ increases to $j+1$, the ratio of the new and old
unnormalized weights on the old support is $(j+1)/(j+1-i)$, which increases in
$i$, and a new positive weight appears at $i=j+1$. So these weights shift
strictly toward larger indices, and the mean of the strictly decreasing sequence
decreases strictly. To check this last step, cross-multiply the 2 means and pair
the terms with indices $i<\ell$. Hence $b_{j+1}/b_j=1+c_j/b_j$ decreases strictly
with $j$, which is the strict log-concavity.

It follows that $B_{a+1}(y)B_{b-1}(y)>B_a(y)B_b(y)$ whenever $b\ge a+2$. Apply
this for $y=vt>0$ in~\eqref{eq:simplex-balance-integral} and integrate. Repeated
transfers give the conclusion.
\end{proof}

\begin{proof}[Proof of Corollary~\ref{cor:simplex-modes}]
For large $N$, $T\le CL$ forces $p>1/d$. So by Theorem~\ref{thm:simplex-balance}
every global mode is a vertex or a balanced bin, and all balanced bins with the
same number of positive coordinates have the same probability. Fix $0<c<1$. The
ratio $C_{N,k}/V_N=S_n(u)$ increases with $u$, and $u$ increases with $T$. So for
$T\le cL$ it is at most its value at $T=cL$, which is
$\exp(-(k-1)(1-c)L+O(\log L))\to0$ by Theorem~\ref{thm:simplex-window}(i). There
the vertices are the only modes, and $p>p_{N,d}$. For $cL\le T\le CL$ we use
Theorem~\ref{thm:simplex-window}(i) in the form
$\log(C_{N,k}/V_N)=(k-1)(\Psi_N(T)-a_k)+o(1)$, uniformly. Fix $\delta>0$, with
$\delta<(a_{d-1}-a_d)/2$ when $d\ge3$. This is possible since $a_2>\cdots>a_d$.
If $\Psi_N(T)\le a_d-\delta$, every face ratio is at most
$e^{-(k-1)\delta+o(1)}<1$, so only the vertices are modes. If
$\Psi_N(T)\ge a_d+\delta$, the full center beats the vertices, and it beats a
balanced bin with $k<d$ positive coordinates by
\[
(d-k)(\Psi_N(T)-a_d)+(k-1)(a_k-a_d)+o(1)>0 .
\]
If $|\Psi_N(T)-a_d|<\delta$, a proper face with $k\ge2$ states has
$(k-1)(\Psi_N(T)-a_k)+o(1)\le-(k-1)\delta+o(1)<0$, since $a_k-a_d>2\delta$. So it
loses to the vertices, and the comparison of the vertices with the full centers
is decided by the unique crossing $p_{N,d}$ of Theorem~\ref{thm:simplex-crossing}.
This proves (i) to (iii) for large $N$. The last claim follows
from~\eqref{eq:simplex-root-expansion} and $p_{N,k}=1-(d-1)z_{N,k}/N$, since
$a_2>\cdots>a_d$.
\end{proof}
